\section*{Contributions}
\addcontentsline{toc}{section}{Contributions}
\label{sec:contributions}

Le projet a été réalisé individuellement. Le tableau~\ref{tab:contrib} détaille les tâches effectuées.

\begin{table}[htbp]
\centering
\caption{Répartition des tâches.}
\label{tab:contrib}
\begin{tabularx}{\textwidth}{>{\bfseries}l X}
\toprule
\textbf{Étudiant} & \textbf{Contribution} \\
\midrule
Abderrahmane EL-HOUIJ &
Recherche bibliographique et rédaction de l'état de l'art ; téléchargement, audit et nettoyage du jeu de données (doublons exacts et quasi-doublons) ; découpage stratifié et validation croisée ; implémentation du pipeline en PyTorch (prétraitement, augmentation sur GPU, tête commune, trois stratégies de transfert, CNN from scratch, fusion de caractéristiques et vote souple) ; entraînement des 65 modèles et des 150 têtes de fusion ; évaluation (métriques, matrices de confusion, courbes ROC, robustesse, coût de calcul) ; visualisations Grad-CAM ; rédaction du rapport et du fichier README ; organisation du dépôt Git. \\
\bottomrule
\end{tabularx}
\end{table}
